\section{Introduction}

\subsection{Background: Spatial Planning and Transport Policies in Switzerland}
Spatial planning in Switzerland operates in a layered fashion.
The federal Spatial Planning Act (LAT: Loi sur l'aménagement du territoire) \cite{noauthor_rs_nodate} was introduced in 1979 by the Confederation and apply to each Cantons and Communes.
It defines the overall framework and general objectives such as containing urban sprawl, protecting natural resources, and ensuring efficient land use, but delegates the implementation to the Cantons. Each Canton is responsible for elaborating a "cantonal master plan" (plan directeur cantonal) within the first eight years of the LAT’s introduction and then revising it at least every ten years. The Confederation must review and approve each plan and its revisions.  

These plans translate federal principles into territorial strategies that account for local geography and specific needs. Once approved by the Confederation, they become binding for public authorities but remain indicative for private actors, shaping how subsequent zoning and land-use decisions are made. The Cantons delegate the responsibility for implementing the measures to the Communes, with the exception of cantonal infrastructure such as cantonal roads, which remain under cantonal responsibility.  

The Communes achieve the goals either by adopting local regulations (règlement communal) or by directly implementing measures, for instance by creating new bike lanes. Some Communes also prepare a "local master plan" (plan d’affectation communal), which materializes these strategies by determining the precise allocation of land and the regulations governing its use.  

This nested structure creates offer flexibility to adapt the general goal as close as possible to the local reality but it require a fair amount of effort to ensure the efficient communication between all the actors. Over time, revisions of the LAT, most notably in 2014 \cite{ssemblee_federale_de_la_confederation_suisse_ro_2014}, shifted the discourse from expansion toward densification, sustainability, and the coordination of land use with transport systems. The vast majority of early plans exist only as paper archives available only trough physical consultation, limiting historical analysis, but comparative studies between cantonal documents remain meaningful. Each plan reveals how spatial and transport policies adapt a shared legal foundation to diverse territorial realities and political cultures.

\newpage

\subsection{Research Motivation: Linking Travel Behavior, Planning, and Text Analysis}

Effective policy design requires a systemic perspective. Interventions that optimize without taking into account the system as a whole often if not always generate unintended consequences\cite{merton_unanticipated_1936}. The problem is actually very similar to what arise in artificial intelligence when systems end up acting in unanticipated ways due to failure (among other factor) of the cost function to capture the entirety of the human intent with the only goal of minimize the said cost function \cite{maier_take_2025}.

his principle directly applies to transportation policy. One of its main goals is to reduce environmental harm by lowering greenhouse gas emissions, which are the primary cause of human-induced climate change. The transport sector plays a key role in this effort. In Europe, it is the only major sector where emissions have not decreased, and it remains the second-largest emitter after energy production \cite{noauthor_co2_2019}. Decarbonizing transport is therefore essential for mitigating climate change. However, greenhouse gases are not the only negative externality. Transport also generates tire-related microplastics, fine particles, noise pollution, accident risks, soil sealing, and barriers to animal movement. These impacts must be balanced with other objectives derived from the right to mobility \cite{nations_universal_nodate}, such as equity and inclusivity, economic efficiency, and safety. Ultimately, transport policy must navigate a tradeoff between protecting the environment and its inhabitants from harmful externalities and ensuring the human need for mobility.

A transition to active mobility and public transport offers the most equitable path forward. Private vehicles are highly inefficient in their use of urban space. They provide primary benefits to a minority of users. This often occurs at a significant cost to the broader community \cite{litman_evaluating_2007}. Prioritizing sustainable modes is a matter of fairness and justice \cite{guo_case_nodate}.

However, the existence of sustainable infrastructure does not guarantee its use. Research, such as the work of J. Grandvillemin, demonstrates this clearly. Modal choice is not a purely rational decision. It is shaped by intangible factors like personal habits, preferences, and perceived access. The sociological concept of motility effectively captures this capacity for mobility \cite{kaufmann_motility_2004}. Understanding these human factors is crucial for designing effective systems.

This leads to the core research question of this document. To what degree are these human dimensions of mobility integrated into Swiss spatial planning? This study investigates whether cantonal master plans consider the factors that truly shape travel behavior. It aims to determine if planning discourse addresses how people experience time and choose between mobility options. The ultimate goal is to assess if Swiss policy aligns physical infrastructure with the human elements of mobility.
\subsection{Research Questions}

To what extent do Swiss territorial / spatial planning documents acknowledge the distinction between perceived time vs measured/travel-time objectively modeled, in relation to modal choice?

How strongly do Swiss planning documents incorporate human / behavioral / motility-based rationales, as opposed to purely economic / efficiency rationales, in shaping transport and spatial policy?
